\begin{frame}{실무 노트: 과적합 감시}
  \begin{itemize}
    \item \kb{3$\sim$4배의 파라미터는 작은 데이터에서 과적합으로
      돌아온다}
    \item \textbf{train $\gg$ val gap}가 \textbf{시퀀스 길이에
      비례해 벌어진다면}:
    \begin{itemize}
      \item 모델 크기를 \textbf{줄이}거나(GRU $\to$ RNN으로 한
        칸)
      \item \textbf{dropout / early stopping}을 더 쓰는 쪽이
      \textbf{게이트를 더 붙이는 것보다 먼저} 시도해야 함
    \end{itemize}
  \end{itemize}
  \vskip0.6em
  \begin{block}{자주 하는 실수: ``소실을 완화했다''고 긴 시퀀스를
  \textbf{무조건 RNN으로}}
    게이트가 소실을 \textbf{느리게} 한다고 해서 \textbf{아무 길이든}
    LSTM/GRU가 ``기억한다''고 믿는 것은 실수. 두 가지 함정:
    \begin{itemize}
      \item \textbf{아직도 지수다, 다만 밑이 1에 가까워졌을 뿐} ---
        수천 토큰 이상이면 \textbf{LSTM/GRU도} 먼 과거를 서서히
        잃어감. ``완화이지 \textbf{해결이 아니다}'' --- 길이가
        극단적일 때는 \textbf{Attention}으로 가는 게 정답
      \item \textbf{과적합} --- 수천 문장 이하의 작은 문제에서
        ``시퀀스가 짧아 소실이 안 생길 것 같은데'' LSTM을 올리면,
        소실은 안 생기고 \textbf{과적합만} 생김 --- ``소실
        증상''(훈련 중에도 긴 시퀀스 정확도가 안 옴)과
        ``과적합 증상''(train $\gg$ val)을 구분할 줄 알아야
      한다
    \end{itemize}
  \end{block}
\end{frame}
